\documentclass{article}
\usepackage[spanish]{babel}
\usepackage[utf8]{inputenc}
\usepackage{graphicx}
\usepackage{amsmath}
\usepackage[a4paper, total={6in, 8in}]{geometry}

\title{Fase 1. Comprensión del problema}

\author{
    Diana Laura Hernández Almeida \\
    Gonzalo Makenly Higuera Inzunza \\
    José Alfredo López Torres \\
    Juan Pablo De La Peña Gonzalez \\
    Sarah Camila Guzmán Fierro
}

\date{Marzo 2026}

\begin{document}

\maketitle

\section*{Introducción}

\textbf{Objetivo:} Entender qué es la calidad del aire, cómo se mide y por qué es un problema relevante.

\section*{Actividades}

\subsection{Investigación sobre contaminantes atmosféricos (PM${2.5}$, PM${10}$, O$_3$, NO$_2$, SO$_2$, CO).} 
    \begin{itemize}
        \item PM$_{2.5}$: Partículas finas con diámetro menor a 2.5 micrómetros, generadas principalmente por procesos de combustión (vehículos, industria). Pueden penetrar profundamente en los pulmones y causar problemas respiratorios y cardiovasculares.
        \begin{itemize}
            \item Valores típicos: ciudad $10$--$35\,\mu g/m^3$, campo $5$--$15\,\mu g/m^3$.
            \item Valores peligrosos: $>55\,\mu g/m^3$.
        \end{itemize}
        
        \item PM$_{10}$: Partículas con diámetro menor a 10 micrómetros, originadas por polvo, construcción y actividades industriales. Afectan principalmente las vías respiratorias superiores.
        \begin{itemize}
            \item Valores típicos: ciudad $20$--$60\,\mu g/m^3$, campo $10$--$30\,\mu g/m^3$.
            \item Valores peligrosos: $>150\,\mu g/m^3$.
        \end{itemize}
        
        \item O$_3$ (ozono troposférico): Contaminante secundario formado por reacciones entre óxidos de nitrógeno y compuestos orgánicos volátiles en presencia de luz solar. Puede irritar las vías respiratorias y afectar cultivos.
        \begin{itemize}
            \item Valores típicos: ciudad $40$--$80$ ppb, campo $30$--$70$ ppb.
            \item Valores peligrosos: $>100$ ppb.
        \end{itemize}
        
        \item NO$_2$ (dióxido de nitrógeno): Gas generado principalmente por la combustión en vehículos y plantas industriales. Contribuye a la formación de smog y lluvia ácida, además de afectar el sistema respiratorio.
        \begin{itemize}
            \item Valores típicos: ciudad $20$--$60$ ppb, campo $5$--$20$ ppb.
            \item Valores peligrosos: $>100$ ppb.
        \end{itemize}
        
        \item SO$_2$ (dióxido de azufre): Emitido por la quema de combustibles fósiles con contenido de azufre, como carbón y petróleo. Puede causar irritación respiratoria y contribuir a la lluvia ácida.
        \begin{itemize}
            \item Valores típicos: ciudad $5$--$20$ ppb, campo $1$--$5$ ppb.
            \item Valores peligrosos: $>75$ ppb.
        \end{itemize}
        
        \item CO (monóxido de carbono): Gas incoloro e inodoro producido por combustión incompleta. Disminuye la capacidad de la sangre para transportar oxígeno, siendo altamente peligroso en concentraciones elevadas.
        \begin{itemize}
            \item Valores típicos: ciudad $0.5$--$5$ ppm, campo $0.1$--$1$ ppm.
            \item Valores peligrosos: $>9$ ppm (exposición prolongada).
        \end{itemize}    
    \end{itemize}

\subsection{Análisis del impacto en salud y medio ambiente.}
    La contaminación del aire representa el noveno factor de riesgo de muerte y discapacidad en México, causando aproximadamente 4.2 millones de muertes prematuras 
    anuales en el mundo. En 2015, reducir las concentraciones de PM2.5 al límite recomendado por la OMS habría evitado hasta 14,666 muertes y más de 150,000 años de 
    vida ajustados por discapacidad en el país. Algunas de las afectaciones en salud incluyen enfermedades respiratorias y cardiovasculares, disfunciones del sistema 
    nervioso central y reproductivo o cáncer, entre otros.
    \begin{table}[h]
        \centering
        \caption{Valores límite en Normas Oficiales Mexicanas y valores guía de la OMS}
        \begin{tabular}{llccccc}
            \hline
             & \textbf{Métrica Exposición} & \textbf{Valor OMS} & \textbf{Valor NOM} & \textbf{Año 1°} & \textbf{Año 3°} & \textbf{Año 5°} \\
            \hline
            PM$_{10}$ & Promedio 24-hr. & 45 & 75 & 70 & 60 & 50 \\
                      & Anual            & 15 & 40 & 36 & 28 & 20 \\
            \hline
            PM$_{2.5}$ & Promedio 24-hr. & 15 & 45 & 41 & 33 & 25 \\
                       & Anual            &  5 & 12 & 10 & 10 & 10 \\
            \hline
            O$_3$ & Máximo 1-hr.         & ---   & 0.095 & 0.090 & 0.090 & 0.090 \\
                  & Promedio móvil 8-hr. & 0.051 & 0.070 & 0.065 & 0.060 & 0.051 \\
            \hline
            NO$_2$ & Máximo 1-hr.    & 0.106 & 0.21  & 0.106 & ---   & ---   \\
                   & Promedio 24-hr. & 0.013 & ---   & ---   & ---   & ---   \\
                   & Anual           & 0.005 & ---   & 0.021 & ---   & ---   \\
            \hline
            SO$_2$ & Máximo 1-hr. & ---   & 0.075 & --- & --- & --- \\
                   & 24 h.         & 0.015 & 0.040 & --- & --- & --- \\
            \hline
        \end{tabular}
        \small\textit{Fuente: INECC / INSP. Unidades: $\mu$g/m$^3$ para PM; ppm para O$_3$, NO$_2$ y SO$_2$.}
    \end{table}
    Los principales contaminantes, partículas PM10 y PM2.5, ozono ($\text{O}_3$), dióxido de nitrógeno ($\text{NO}_2$) y dióxido de azufre ($\text{SO}_2$); afectan 
    el sistema respiratorio, cardiovascular y nervioso, con mayor impacto en niños, adultos mayores y mujeres embarazadas. Las partículas más pequeñas (< 2.5 micras) 
    penetran hasta los alvéolos pulmonares e incluso el torrente sanguíneo. En 2019, ninguna de las 25 ciudades mexicanas evaluadas cumplió con los límites normados 
    de PM2.5, y Guadalajara y Monterrey registraron las concentraciones más elevadas de PM10 y $\text{O}_3$.


\subsection{Sesión de interacción con estudiantes de Ingeniería Química para discutir:}
    \paragraph{Smog Fotoquímico}
    El smog fotoquímico se forma a partir de una serie de reacciones químicas en la atmósfera, principalmente impulsadas por la radiación solar. Uno de los procesos clave es la fotólisis del dióxido de nitrógeno:
    
    \begin{align*}
    \text{NO}_2 + \text{luz solar} &\rightarrow \text{NO} + \text{O} \\
    \text{O} + \text{O}_2 &\rightarrow \text{O}_3
    \end{align*}
    
    En este proceso, el NO$_2$ se descompone por efecto de la radiación ultravioleta, liberando un átomo de oxígeno altamente reactivo que posteriormente forma ozono (O$_3$), un contaminante secundario.
    
    Por otro lado, el monóxido de carbono (CO) participa indirectamente en la formación de ozono mediante reacciones con radicales libres:
    
    \begin{align*}
    \text{CO} + \cdot\text{OH} &\rightarrow \text{CO}_2 + \cdot\text{H} \\
    \cdot\text{H} + \text{O}_2 &\rightarrow \cdot\text{HO}_2 \\
    \cdot\text{HO}_2 + \text{NO} &\rightarrow \text{NO}_2 + \cdot\text{OH}
    \end{align*}
    
    Este conjunto de reacciones genera un ciclo químico en el que el NO$_2$ se regenera, permitiendo la producción continua de ozono en la atmósfera.
    
    \paragraph{Formación de lluvia ácida a partir de SO$_2$}
    
    El SO$_2$ se disuelve en agua atmosférica formando H$_2$SO$_3$, que posteriormente se oxida a H$_2$SO$_4$. Este proceso ocurre principalmente en condiciones de alta humedad y presencia de oxidantes.

\section*{Referencias}

\begin{itemize}

\item INECC, S. Informe nacional de calidad del aire 2015. 2016. México.

\item Institute for Health Metrics and Evaluation (IHME). Mexico profile. Seattle, WA: IHME, University of Washington. 2018. Available from http://www.healthdata.org/Mexico (Accessed august 2021). 

\item Manisalidis I, Stavropoulou E, Stavropoulos A, Bezirtzoglou E. Environmental and health impacts of air pollution: a review. Frontiers in public health. 2020;8(14). 
\\ https://doi.org/10.3389/fpubh.2020.00014

\item ONU-Habitat México. Índice de prosperidad urbana en la República Mexicana. Reporte Nacional de Tendencias de la Prosperidad Urbana en México. México: ONU-Habitat, 2016.

\item Trejo-González AG, Riojas-Rodriguez H, Texcalac-Sangrador JL, Guerrero-López CM, Cervantes-Martínez K, Hurtado-Díaz M, Zuñiga-Bello PE. Quantifying health impacts and economic costs of PM 2.5 exposure in Mexican cities of the National Urban System. International journal of public health. 2019;64(4):561-72. https://doi.org/10.1007/s00038-019-01216-1

\item Partículas suspendidas PM10 Y PM2.5 dañan salud y medio ... (n.d.). \\
https://www.gob.mx/semarnat/articulos/particulas-suspendidas-pm10-y-pm2-5-danan-salud-y-medio-ambiente
\end{itemize}

\end{document}